\begin{center}
\textsc{Gauss} an \textsc{Olbers}.\quad Göttingen, 19.~Februar 1825.
\end{center}

\par\noindent\hspace*{\parindent}Ich schicke Ihnen mein vollständiges Höhenverzeichnis[\textsuperscript{*})].
Bei Hamburg ist eine kleine Veränderung von \(\pm\,2{,}7\) Fuß vorgenommen, da \textsc{Schumacher} mir
den Höhenunterschied zwischen dem Knopf und den Fenstern des Kabinetts,
den ich früher zu 47,9 Fuß nach \textsc{Benzenbergs} Kupfer[stich] angenommen
hatte, nach \textsc{Sonnins} Kupfer[stich] zu 53,3 angibt; ich habe also einstweilen den
Knopf (meinen Zielpunkt) um 2,7 höher, die Fenster 2,7 [Fuß] tiefer gesetzt
als vorher.
Dadurch wird dann auch \textsc{Schumachers} Barometer 2,7 [Fuß] tiefer
als vorher, also gegen Göttingen \(-\,357{,}8\) [Fuß] und gegen Ihr Barometer \(+\,70{,}4\)
[Fuß] \(= 11{,}73\) Toisen.
Der Unterschied mit Ihrer Barometerbestimmung ist
also jetzt ganz unbedeutend.
Doch wird dies noch etwas modifiziert werden, da
ich auch auf einigen Stationen den Fußboden der Laterne zum Zielpunkt
gebraucht habe, dessen Tiefe unter dem Knopfe mir \textsc{Schumacher} nicht mitgeteilt hat,
und überhaupt sollte wohl die relative Höhe der drei Punkte
ordentlich trigonometrisch gemessen werden (aus einem nahen Standpunkte).

\par\smallskip\noindent\rule{.24\linewidth}{.4pt}\par\smallskip
{\footnotesize
\noindent\hspace*{2.75em}[\textsuperscript{*}) In Gauss’ Nachlass ist das Original dieses Verzeichnisses vorhanden; demselben ist das Folgende zugefügt:]\par
\noindent\hspace*{2.75em}«Die Höhen bei den Hauptpunkten beziehen sich, mit Ausnahme von Göttingen, Lüneburg und Hamburg,
wo das Nähere besonders bemerkt ist, allemal auf diejenige Fläche, auf welcher der Repetitionskreis gestanden hat.
Bei den meisten Plätzen ist dies ein aufgemauertes steinernes Postament von 3\(\tfrac{1}{2}\) Fuß Höhe über der Erde,
zuweilen ein paar Zoll mehr, zuweilen weniger.
Nur beim Meridianzeichen beträgt diese Höhe etwas mehr, nämlich 5 Fuß«.\par
}

\clearpage
\noindent\null\hfill{\small\textbf{ZUR HANNOVERSCHEN TRIANGULATION.}}\hfill\makebox[0pt][r]{373}\par\medskip
\begin{center}
\renewcommand{\arraystretch}{1.2}
Relative Höhen.\par\medskip
\begin{tabular}{|r|p{.70\linewidth}|c@{\hspace{.35em}}r|}
\hline
   & & \multicolumn{2}{c|}{Pariser Fuß} \\
 1 & Göttingen, Sternwarte, Fußboden & & \(0\phantom{{,}0}\) \\
 2 & Nördl.\ Meridianzeichen (Oberfl.\ des Postaments) & \(+\) & \(201{,}0\) \\
 3 & Hohehagen (Postamentsoberfläche) & \(+\) & \(1072{,}6\) \\
 4 & Hils (P.) & \(+\) & \(841{,}2\) \\
 5 & Brockenhaus (Marmortisch auf dem Turm) & \(+\) & \(3061{,}7\) \\
 6 & Lichtenberg (P.) & \(+\) & \(274{,}1\) \\
 7 & Deister, Calenberg (P.) & \(+\) & \(468{,}2\) \\
 8 & Garssen (P.) & \(-\) & \(242{,}5\) \\
 9 & Falkenberg (P.) & \(-\) & \(15{,}5\) \\
10 & Scharnhorst (P.) & \(-\) & \(190{,}6\) \\
11 & Breithorn (P.) & \(-\) & \(109{,}0\) \\
12 & Hauselberg (P.) & \(-\) & \(109{,}2\) \\
13 & Wulfsode (P.) & \(-\) & \(158{,}5\) \\
14 & Timpenberg (P.) & \(-\) & \(120{,}1\) \\
15 & Nindorf (P.) & \(-\) & \(120{,}0\) \\
16 & Lüneburg, Michaelis, Fußboden der Laterne & \(-\) & \(243{,}2\) \\
17 & \hspace*{1.5em}\makebox[4.45em][l]{---}\makebox[3.15em][l]{---}Knopf & \(-\) & \(178{,}6\) \\
18 & \hspace*{1.5em}\makebox[4.45em][l]{---}\makebox[3.15em][l]{---}Spitze & \(-\) & \(168{,}3\) \\
19 & \hspace*{1.5em}\makebox[3.15em][l]{---}Johannis, Knopf & \(-\) & \(109{,}5\) \\
20 & \hspace*{1.5em}\makebox[4.45em][l]{---}\makebox[2.95em][l]{---}Spitze & \(-\) & \(101{,}4\) \\
21 & \hspace*{1.5em}\makebox[3.15em][l]{---}Nicolai, Knopf & \(-\) & \(199{,}3\) \\
22 & \hspace*{1.5em}\makebox[4.3em][l]{---}\makebox[2.4em][l]{---}Spitze & \(-\) & \(188{,}4\) \\
23 & \hspace*{1.5em}\makebox[2.95em][l]{---}Lamberti, Knopf & \(-\) & \(215{,}7\) \\
24 & \hspace*{1.5em}\makebox[4.45em][l]{---}\makebox[2.95em][l]{---}Spitze & \(-\) & \(207{,}7\) \\
25 & \hspace*{1.5em}\makebox[2.95em][l]{---}Höchste Stelle des Kalkberges & \(-\) & \(291{,}6\) \\
26 & \hspace*{1.5em}\makebox[2.95em][l]{---}Platz nahe vor dem N.W.~Tor & \(-\) & \(393{,}1\) \\
27 & Lüne, Klosterturm, Knopf & \(-\) & \(331{,}2\) \\
28 & \hspace*{.55em}\makebox[4.45em][l]{---}\makebox[3.9em][l]{---}Spitze & \(-\) & \(325{,}0\) \\
29 & Hamburg, Michaelisthurm, Knopf & \(-\) & \(31{,}2\) \\
\end{tabular}
\end{center}

\par\clearpage
\noindent\makebox[0pt][l]{374}\hfill{\small\textbf{BRIEFWECHSEL MIT OLBERS.}}\hfill\null\par\medskip

\begin{center}
\renewcommand{\arraystretch}{1.2}
\begin{tabular}{|r|p{.70\linewidth}|c@{\hspace{.35em}}r|}
\hline
 & & \multicolumn{2}{c|}{Pariser Fuß} \\
30 & Hamburg, Michaelisthurm, Fenster des ob.\ Kab. & \(-\) & \(84{,}5\) \\
31 & Wilsede (P.) & \(+\) & \(45{,}8\) \\
32 & Elmhorst (P.) & \(-\) & \(202{,}9\) \\
33 & Litberg (P.) & \(-\) & \(278{,}1\) \\
34 & Bullerberg (P.) & \(-\) & \(315{,}7\) \\
35 & Bottel (P.) & \(-\) & \(318{,}5\) \\
36 & Steinberg (P.) & \(-\) & \(257{,}9\) \\
37 & Brüttendorf (P.) & \(-\) & \(326{,}3\) \\
38 & Zeven, Kirchturm, Fußboden der Laterne & \(-\) & \(347{,}0\) \\
39 & \hspace*{.6em}\makebox[5.05em][l]{---}\makebox[4.1em][l]{---}Knopf & \(-\) & \(323{,}8\) \\
40 & \hspace*{.6em}\makebox[5.05em][l]{---}\makebox[4.1em][l]{---}Spitze & \(-\) & \(316{,}8\) \\
41 & \hspace*{.6em}\makebox[3em][l]{---}Garten beim Posthause, nahe der Aue & \(-\) & \(430{,}5\) \\
42 & \hspace*{.6em}\makebox[3em][l]{---}Platz auf dem Felde beim südlichen Eingang des\newline\hspace*{4.65em}Dorfs & \raisebox{-\normalbaselineskip}[0pt][0pt]{\(-\)} & \raisebox{-\normalbaselineskip}[0pt][0pt]{\(419{,}4\)} \\
43 & Bremen, Ansgarius, Fußboden der Laterne & \(-\) & \(227{,}7\) \\
44 & \hspace*{1.2em}\makebox[5em][l]{---}\makebox[3.4em][l]{---}Knopf & \(-\) & \(172{,}1\) \\
45 & \hspace*{1.2em}\makebox[5em][l]{---}\makebox[3.4em][l]{---}Spitze & \(-\) & \(156{,}6\) \\
46 & \hspace*{1.2em}\makebox[3em][l]{---}Liebenfrauen, Knopf & \(-\) & \(215{,}4\) \\
47 & \hspace*{1.2em}\makebox[3em][l]{---}Dom, Knopf & \(-\) & \(244{,}2\) \\
48 & \hspace*{1.2em}\makebox[3.6em][l]{---}\makebox[2.15em][l]{---}Spitze & \(-\) & \(232{,}4\) \\
49 & \hspace*{1.2em}\makebox[3em][l]{---}Martini, Knopf & \(-\) & \(281{,}7\) \\
50 & \hspace*{1.2em}\makebox[3em][l]{---}Gymnasium, Knopf & \(-\) & \(315{,}9\) \\
51 & \hspace*{1.2em}\makebox[3em][l]{---}Straßenpflaster am Domhof unweit des alten Mu-\newline\hspace*{5.25em}seums & \raisebox{-\normalbaselineskip}[0pt][0pt]{\(-\)} & \raisebox{-\normalbaselineskip}[0pt][0pt]{\(453{,}8\)} \\
52 & \hspace*{1.2em}\makebox[3em][l]{---}Windmühlenberg am Herdner Tor & \(-\) & \(443{,}3\) \\
53 & \hspace*{1.2em}\makebox[3em][l]{---}Garten beim Hause des Hrn.\ Dr.\ \textsc{Olbers} & \(-\) & \(452{,}7\) \\
54 & \hspace*{1.2em}\makebox[3em][l]{---}Fußboden in Hrn.\ Dr.\ \textsc{Olbers} Observ. & \(-\) & \(430{,}7\) \\
55 & \hspace*{1.2em}\makebox[3em][l]{---}Barometergefäß, daselbst & \(-\) & \(428{,}2\) \\
56 & Brillit (P.) & \(-\) & \(342{,}2\) \\
57 & Garlster Haide (P.) & \(-\) & \(328{,}7\) \\
\end{tabular}
\end{center}
